\uuid{Ty8y}
\titre{Reconnaître les paramètres d'une densité normale}
\auteur{Maxime Nguyen}
\organisation{AMSCC}
\datecreate{2026-11-27}
\niveau{L2}
\module{Probabilité}
\chapitre{Probabilité continue}
\sousChapitre{Loi normale}
\difficulte{1}
\theme{loi normale, densité, paramètres}

\contenu{
\begin{question}
\qtype{onlyone}
On note $\mathcal{N}(\mu,\sigma)$ la loi normale d'espérance $\mu$ et d'écart-type $\sigma$ ; sa densité est $x \mapsto \frac{1}{\sigma\sqrt{2\pi}}\,e^{-\frac{(x-\mu)^2}{2\sigma^2}}$. Une variable $X$ a pour densité
$$f(x) = \frac{1}{3\sqrt{2\pi}}\,e^{-\frac{(x+1)^2}{18}}.$$
La loi de $X$ est :
\begin{answers}
\good{$\mathcal{N}(-1, 3)$}
\bad{$\mathcal{N}(1, 3)$}
\bad{$\mathcal{N}(-1, 9)$}
\bad{$\mathcal{N}(-1, 18)$}
\end{answers}
\begin{explanations}
On lit $x - \mu = x + 1$, donc $\mu = -1$, et
$2\sigma^2 = 18$, donc $\sigma^2 = 9$ et $\sigma = 3$, ce que confirme le
facteur $\frac{1}{3\sqrt{2\pi}}$.

B : le signe de $\mu$ est mal lu ($x + 1 = x - (-1)$).

C donne la variance au lieu de l'écart-type, D le dénominateur $2\sigma^2$.

Contre-sens visés : lire $\mu$ avec le mauvais signe ; confondre $\sigma$,
$\sigma^2$ et $2\sigma^2$.
\end{explanations}
\end{question}
}
